% =======================================================
% ВАРИАТИВНЫЕ ШАБЛОНЫ и МЕТАПРОГРАММИРОВАНИЕ
% =======================================================

\section{Вариативные шаблоны (variadic templates)}

% -------------------------------------------------------
\subsection{Что это}

\textbf{Вариативный шаблон} принимает \textbf{произвольное число} параметров
любых типов. Это позволило реализовать в библиотеке \texttt{printf}-подобные
функции, \texttt{std::tuple}, \texttt{std::make\_unique} и др.

\begin{lstlisting}
template<typename... Args>      // Args -- пакет параметров типов
void print(Args... args);      // args -- пакет аргументов-значений
print(1, 2.5, "hello");        // Args = {int, double, const char*}
\end{lstlisting}

\subsubsection{Пакеты параметров (parameter packs)}

\begin{itemize}[leftmargin=2em]
    \item \texttt{typename... Args} --- пакет шаблонных параметров.
    \item \texttt{Args... args} --- пакет функциональных параметров.
    \item \texttt{args...} --- \textbf{развёртка (expansion)} пакета.
    \item \texttt{sizeof...(Args)} --- число элементов в пакете.
\end{itemize}

\subsection{Рекурсивная обработка пакета (классический способ)}

До C++17 пакет обрабатывали рекурсией: «голова + хвост» с базовым случаем:

\begin{lstlisting}
void print() {}   // базовый случай: пустой пакет

template<typename Head, typename... Tail>
void print(Head h, Tail... tail) {
    std::cout << h;                 // обработали первый
    if (sizeof...(tail) > 0) std::cout << ", ";
    print(tail...);                 // рекурсия по остатку
}
print(1, 2, 3);   // 1, 2, 3
\end{lstlisting}

\subsection{Fold-выражения (C++17)}

\textbf{Свёртки (fold expressions)} применяют оператор ко всему пакету без
рекурсии:

\begin{lstlisting}
template<typename... Args>
auto sum(Args... args) {
    return (args + ...);      // унарная правая свёртка: a1 + (a2 + (a3 + ...))
}
sum(1, 2, 3, 4);   // 10
\end{lstlisting}

Четыре формы свёрток:

\begin{center}
\begin{tabular}{ll}
\toprule
\textbf{Форма} & \textbf{Раскрытие} \\
\midrule
\texttt{(pack op ...)}          & унарная правая: $e_1\ op\ (\dots op\ e_n)$ \\
\texttt{(... op pack)}          & унарная левая:  $((e_1\ op\ e_2)\ op\ \dots)$ \\
\texttt{(pack op ... op init)}  & бинарная правая (с начальным значением) \\
\texttt{(init op ... op pack)}  & бинарная левая \\
\bottomrule
\end{tabular}
\end{center}

\begin{lstlisting}
template<typename... Args>
void print(Args... args) {
    ((std::cout << args << ' '), ...);   // свёртка по запятой
}
template<typename... Args>
bool all_true(Args... args) {
    return (args && ...);                // все ли истинны
}
\end{lstlisting}

\subsection{Perfect forwarding в вариативных шаблонах}

Универсальная фабрика передаёт аргументы дальше \textbf{без потери} категории
(lvalue/rvalue) через \texttt{std::forward}:

\begin{lstlisting}
template<typename T, typename... Args>
std::unique_ptr<T> make(Args&&... args) {
    return std::unique_ptr<T>(new T(std::forward<Args>(args)...));
}
\end{lstlisting}

\texttt{Args\&\&...} --- пакет \textbf{forwarding-ссылок};
\texttt{std::forward<Args>(args)...} разворачивает пакет, сохраняя категорию
каждого аргумента.

\subsection{Пример: типобезопасный кортеж и доступ по индексу}

\begin{lstlisting}
template<typename... Ts>
struct Tuple {};   // упрощённо; в стандарте -- std::tuple

std::tuple<int, std::string, double> t{1, "hi", 3.14};
auto x = std::get<1>(t);              // "hi"
auto [a, b, c] = t;                   // structured bindings
\end{lstlisting}

\subsection{Развёртка в разных контекстах}

\begin{lstlisting}
template<typename... Args>
void demo(Args... args) {
    f(args...);                  // f(a1, a2, a3)
    g(h(args)...);               // g(h(a1), h(a2), h(a3))
    int arr[] = {args...};       // инициализация массива
    std::vector v{args...};
}
\end{lstlisting}


% =======================================================
\section{Метапрограммирование: constexpr, метафункции, SFINAE, концепты}

% -------------------------------------------------------
\subsection{Что такое метапрограммирование}

\textbf{Шаблонное метапрограммирование (TMP)} --- вычисления, выполняемые
\textbf{на этапе компиляции}. Результат --- типы и константы, «зашитые» в
программу, без затрат в runtime.

\subsection{\texttt{constexpr} --- вычисления на этапе компиляции}

\texttt{constexpr}-функция может быть вычислена компилятором, если её аргументы
известны на этапе компиляции:

\begin{lstlisting}
constexpr int factorial(int n) {
    return (n <= 1) ? 1 : n * factorial(n - 1);
}
constexpr int f5 = factorial(5);   // 120, вычислено компилятором
int arr[factorial(4)];             // 24 -- можно как размер массива
\end{lstlisting}

\begin{itemize}[leftmargin=2em]
    \item \texttt{const} --- неизменяемость (может вычисляться в runtime).
    \item \texttt{constexpr} --- \textbf{может} вычисляться на этапе компиляции.
    \item \texttt{consteval} (C++20) --- \textbf{обязан} вычисляться на этапе
          компиляции (immediate function).
    \item \texttt{constinit} (C++20) --- гарантирует статическую инициализацию.
\end{itemize}

\subsubsection{\texttt{if constexpr} (C++17)}

Ветвление на этапе компиляции: «неподходящая» ветка \textbf{не компилируется}:

\begin{lstlisting}
template<typename T>
auto to_string(T value) {
    if constexpr (std::is_arithmetic_v<T>)
        return std::to_string(value);    // только для чисел
    else
        return value;                    // для остального
}
\end{lstlisting}

\subsection{Метафункции и type traits}

\textbf{Метафункция} --- шаблон, «возвращающий» тип или значение через вложенный
\texttt{::type} или \texttt{::value}. Стандартные метафункции живут в
\texttt{<type\_traits>}.

\begin{lstlisting}
#include <type_traits>
std::is_integral<int>::value;        // true
std::is_integral_v<int>;             // true (суффикс _v -- значение)
std::remove_const<const int>::type;  // int
std::remove_const_t<const int>;      // int (суффикс _t -- тип)
std::is_same_v<int, int32_t>;        // true/false в зависимости от платформы
\end{lstlisting}

\subsubsection{Своя метафункция}

\begin{lstlisting}
// "вернуть" тип: убрать указатель
template<typename T> struct remove_ptr        { using type = T; };
template<typename T> struct remove_ptr<T*>    { using type = T; };  // спец-ция
template<typename T> using remove_ptr_t = typename remove_ptr<T>::type;

// "вернуть" значение: длина массива на этапе компиляции
template<typename T, size_t N>
constexpr size_t array_size(T(&)[N]) { return N; }
\end{lstlisting}

\subsection{SFINAE}

\textbf{SFINAE (Substitution Failure Is Not An Error)}: если при подстановке
шаблонных аргументов сигнатура оказывается некорректной, эта перегрузка
\textbf{молча исключается} из набора кандидатов, а не вызывает ошибку.

\begin{lstlisting}
#include <type_traits>
// версия только для целочисленных T
template<typename T,
         std::enable_if_t<std::is_integral_v<T>, int> = 0>
void process(T x) { std::cout << "integral\n"; }

// версия только для вещественных T
template<typename T,
         std::enable_if_t<std::is_floating_point_v<T>, int> = 0>
void process(T x) { std::cout << "floating\n"; }

process(42);     // "integral"
process(3.14);   // "floating"
\end{lstlisting}

\subsubsection{Идиома detection (проверка наличия члена)}

\begin{lstlisting}
// есть ли у T метод .size()?
template<typename T, typename = void>
struct has_size : std::false_type {};

template<typename T>
struct has_size<T, std::void_t<decltype(std::declval<T>().size())>>
    : std::true_type {};

has_size<std::vector<int>>::value;  // true
has_size<int>::value;               // false
\end{lstlisting}

\subsection{Концепты (C++20)}

\textbf{Концепты} --- именованные ограничения на типы. Они заменяют громоздкий
SFINAE, дают \textbf{читаемые ошибки} и явно документируют требования.

\begin{lstlisting}
#include <concepts>
template<typename T>
concept Numeric = std::is_arithmetic_v<T>;   // своё ограничение

template<Numeric T>          // T обязан быть числом
T square(T x) { return x * x; }

square(5);      // OK
// square("s"); // понятная ошибка: "s" не удовлетворяет Numeric
\end{lstlisting}

\subsubsection{Способы применить ограничение}

\begin{lstlisting}
template<std::integral T> void f(T);            // на параметре шаблона
template<typename T> requires std::integral<T>  // requires-clause
void g(T);
void h(std::integral auto x);                   // сокращённый синтаксис
\end{lstlisting}

\subsubsection{\texttt{requires}-выражение --- задаём свой концепт}

\begin{lstlisting}
template<typename T>
concept Container = requires(T c) {
    c.begin();              // должен быть begin()
    c.end();                // должен быть end()
    c.size();               // должен быть size()
    typename T::value_type; // должен быть вложенный тип value_type
};
\end{lstlisting}

\subsection{Зачем всё это: библиотечный код}

Метапрограммирование --- основа стандартной библиотеки: \texttt{std::vector}
выбирает оптимальные операции по traits типа, алгоритмы диспетчеризуются по
категории итератора, \texttt{std::variant}/\texttt{std::tuple} построены на
вариативных шаблонах. В прикладном коде стоит предпочитать \texttt{if constexpr}
и концепты как наиболее читаемые инструменты.
